%==============================================================================
%  ANNEXE - Complements : donnees, IHM, securite et tests
%  A inserer via %==============================================================================
%  ANNEXE - Complements : donnees, IHM, securite et tests
%  A inserer via %==============================================================================
%  ANNEXE - Complements : donnees, IHM, securite et tests
%  A inserer via %==============================================================================
%  ANNEXE - Complements : donnees, IHM, securite et tests
%  A inserer via \input{annexe_complements} dans cahier_des_charges.tex
%  (charte "Miel & noir" : miel, mielfonce, encadre, colonne C, longtable,
%   booktabs -- deja definis dans le preambule principal)
%==============================================================================

%==============================================================================
\chapter{Annexe - Complements : donnees, IHM, securite et tests}
\label{chap:complements}
%==============================================================================
Cette annexe regroupe des livrables complementaires renforcant le dossier :
le modele logique de donnees, les maquettes de l'interface, la matrice des
droits d'acces, le registre des risques, la conformite a la protection des
donnees et la strategie de tests.

%------------------------------------------------------------------------------
\section{Modele logique de donnees (MLD)}
%------------------------------------------------------------------------------
Le MLD traduit le dictionnaire de donnees (chapitre~5) en tables relationnelles
avec cles primaires, cles etrangeres et contraintes de gestion (\texttt{CHECK}
sur le nombre de hausses, la capacite en cadres et la productivite ; unicite du
couple compartiment/slot). Il distingue la relation d'\emph{emplacement}
(un site heberge des ruches) de la relation de \emph{propriete} (une ferme
possede des ruches), conformement a la regle autorisant un site a accueillir des
ruches de plusieurs fermes.

\begin{figure}[htbp]
\centering
\includegraphics[width=0.82\linewidth]{diagramme_mld}
\caption{Modele logique de donnees relationnel du systeme Z\"umm.}
\label{fig:mld}
\end{figure}

%------------------------------------------------------------------------------
\section{Maquettes de l'interface (wireframes)}
%------------------------------------------------------------------------------
Les maquettes basse-fidelite suivantes fixent l'ergonomie des ecrans cles avant
le developpement. Elles respectent la charte visuelle et la simplicite exigee.

\begin{figure}[htbp]
\centering
\includegraphics[width=0.9\linewidth]{wireframe_calendrier}
\caption{Maquette : calendrier / matrice agents x ruches avec pop-up de visite.}
\label{fig:wf-calendrier}
\end{figure}

\begin{figure}[htbp]
\centering
\includegraphics[width=0.9\linewidth]{wireframe_tdb_production}
\caption{Maquette : tableau de bord de production avec seuil parametrable.}
\label{fig:wf-production}
\end{figure}

\begin{figure}[htbp]
\centering
\includegraphics[width=0.95\linewidth]{wireframe_rapport_visite}
\caption{Maquette : formulaire de rapport de visite (avec mode hors-ligne et photos).}
\label{fig:wf-rapport}
\end{figure}

\begin{figure}[htbp]
\centering
\includegraphics[width=0.85\linewidth]{wireframe_ruche}
\caption{Maquette : fiche d'une ruche et composition (corps, hausses, cadres).}
\label{fig:wf-ruche}
\end{figure}

%------------------------------------------------------------------------------
\section{Matrice des droits d'acces (RBAC)}
%------------------------------------------------------------------------------
La matrice precise, pour chaque fonction, le droit accorde a chaque profil.
\textbf{Legende :} CRUD = creation / lecture / mise a jour / suppression ;
R = lecture ; X = execution ; A = approbation ; --- = aucun acces.

\begin{longtable}{>{\bfseries}p{3.6cm} C{1.3cm} C{1.3cm} C{1.3cm} C{1.3cm} C{1.3cm} C{1.3cm}}
\toprule
\textbf{Fonction} & \textbf{Prop.} & \textbf{Apic.} & \textbf{Sup.} & \textbf{Resp.} & \textbf{Admin} & \textbf{API} \\
\midrule
\endhead
Fermes et sites          & CRUD & R    & R    & R    & CRUD & --- \\
Ruches (corps/cadres)    & CRUD & R    & R    & R    & CRUD & --- \\
Agents                   & R    & ---  & R    & R    & CRUD & --- \\
Planifier une visite     & R    & X    & X    & X    & X    & --- \\
Approuver le planning    & ---  & ---  & A    & ---  & X    & --- \\
Realiser visite/rapport  & R    & X    & X    & ---  & R    & --- \\
Calendrier agents x ruches & R  & R    & R    & R    & R    & --- \\
TB production            & R    & ---  & ---  & R    & R    & --- \\
TB alertes / traiter     & R    & ---  & ---  & X    & R    & --- \\
Synthese et ROI          & R    & ---  & ---  & R    & R    & --- \\
Parametrer seuils/unites & ---  & ---  & ---  & R    & CRUD & --- \\
Internationalisation     & ---  & ---  & ---  & ---  & CRUD & --- \\
Export CSV / TXT         & R    & R    & R    & R    & R    & --- \\
Sauvegarde / restauration & --- & ---  & ---  & ---  & X    & --- \\
API quantite de miel     & ---  & ---  & ---  & ---  & ---  & R \\
\bottomrule
\end{longtable}

\begin{encadre}
Cette matrice doit etre appliquee cote serveur (verification systematique des
droits dans la couche de controle), et non uniquement par le masquage des
elements d'interface.
\end{encadre}

%------------------------------------------------------------------------------
\section{Registre des risques}
%------------------------------------------------------------------------------
\begin{longtable}{>{\bfseries}p{4.6cm} C{1.8cm} C{1.6cm} p{5.4cm}}
\toprule
\textbf{Risque} & \textbf{Probab.} & \textbf{Impact} & \textbf{Mesure de reduction} \\
\midrule
\endhead
Retard de developpement       & Moyenne & Fort      & Prioriser le noyau CRUD ; jalons hebdomadaires (S5 a S12). \\
Portee trop large             & Moyenne & Fort      & Se limiter aux fonctions de priorite haute ; le reste en option. \\
Complexite i18n arabe (RTL)   & Moyenne & Moyen     & Test precoce de bout en bout FR/EN/AR, gestion \texttt{dir=rtl}. \\
Perte de donnees terrain      & Moyenne & Fort      & File de commandes hors-ligne + synchronisation + sauvegarde. \\
Failles applicatives (injection, XSS) & Moyenne & Fort & \texttt{PreparedStatement}, echappement JSP, revue de securite. \\
Fausses alertes (capteurs)    & Moyenne & Moyen     & Seuils parametrables + validation humaine par le rapport. \\
Indisponibilite d'OAuth Google & Faible & Moyen     & Repli d'authentification locale pour la demonstration. \\
Volumetrie des mesures        & Faible  & Moyen     & Archivage et agregation par periode (volet automatise). \\
Usage d'un generateur de code & Faible  & Tres fort & Conception et code faits maison, tracabilite et revue par pairs. \\
\bottomrule
\end{longtable}

%------------------------------------------------------------------------------
\section{Protection des donnees personnelles}
%------------------------------------------------------------------------------
Le systeme manipule des donnees personnelles (noms et contacts des fermiers et
des agents) et des donnees potentiellement sensibles (\textbf{geolocalisation}
precise des sites). Les principes suivants sont retenus, dans l'esprit du RGPD :
\begin{itemize}
  \item \textbf{Minimisation et finalite} : ne collecter que les donnees utiles au suivi apicole.
  \item \textbf{Securite} : chiffrement des echanges (SSL / X.509) et acces authentifies, deja prevus.
  \item \textbf{Controle d'acces} : application stricte de la matrice RBAC ci-dessus.
  \item \textbf{Droits des personnes} : consultation, rectification et \textbf{effacement} des donnees, \textbf{portabilite} assuree par l'export CSV / TXT.
  \item \textbf{Conservation} : duree de retention definie pour les mesures et les visites, au-dela de laquelle les donnees sont archivees ou anonymisees.
  \item \textbf{Tracabilite} : journalisation des acces et des actions sensibles (approbation de planning, cloture de site).
\end{itemize}

%------------------------------------------------------------------------------
\section{Strategie de tests}
%------------------------------------------------------------------------------
La strategie complete le plan de tests fonctionnels (chapitre~10) par les
niveaux unitaire et d'integration. La conception faiblement couplee (interfaces
\texttt{Repository}, services segreges) rend les regles metier testables en
isolation, par injection de doubles (mocks).

\begin{longtable}{>{\bfseries}p{2.8cm} p{4.4cm} p{4.4cm} C{2.2cm}}
\toprule
\textbf{Niveau} & \textbf{Portee} & \textbf{Exemples} & \textbf{Outil} \\
\midrule
\endhead
Unitaire      & Regles de gestion isolees          & Contraintes cadres/hausses, calcul du miel (Composite), ROI, seuils & JUnit 5, Mockito \\
Integration   & Acces aux donnees (Repository/JDBC) & CRUD persiste, requetes de calcul                                   & JUnit + base H2 \\
Systeme       & Cas fonctionnels de bout en bout    & Cas TC01 a TC16 du chapitre~10                                      & JUnit, Selenium \\
Acceptation   & Validation par la demonstration     & Parcours agent, fermier, responsable                                & Soutenance \\
\bottomrule
\end{longtable}

\begin{itemize}
  \item \textbf{Couverture visee} : au moins 70\,\% de la couche metier (regles de gestion), mesuree par \texttt{JaCoCo}.
  \item \textbf{Tracabilite} : chaque exigence est reliee a au moins un cas de test (matrice du chapitre~10), completee par les tests unitaires correspondants.
  \item \textbf{Automatisation} : execution des tests a chaque build (Maven), pre-requis a la livraison.
\end{itemize}

\begin{encadre}
\textbf{Synergie avec la conception :} l'inversion de dependance (DIP) permet de
substituer un \texttt{RucheRepository} en memoire aux implementations JDBC, ce
qui rend les tests rapides, deterministes et independants de la base.
\end{encadre}
 dans cahier_des_charges.tex
%  (charte "Miel & noir" : miel, mielfonce, encadre, colonne C, longtable,
%   booktabs -- deja definis dans le preambule principal)
%==============================================================================

%==============================================================================
\chapter{Annexe - Complements : donnees, IHM, securite et tests}
\label{chap:complements}
%==============================================================================
Cette annexe regroupe des livrables complementaires renforcant le dossier :
le modele logique de donnees, les maquettes de l'interface, la matrice des
droits d'acces, le registre des risques, la conformite a la protection des
donnees et la strategie de tests.

%------------------------------------------------------------------------------
\section{Modele logique de donnees (MLD)}
%------------------------------------------------------------------------------
Le MLD traduit le dictionnaire de donnees (chapitre~5) en tables relationnelles
avec cles primaires, cles etrangeres et contraintes de gestion (\texttt{CHECK}
sur le nombre de hausses, la capacite en cadres et la productivite ; unicite du
couple compartiment/slot). Il distingue la relation d'\emph{emplacement}
(un site heberge des ruches) de la relation de \emph{propriete} (une ferme
possede des ruches), conformement a la regle autorisant un site a accueillir des
ruches de plusieurs fermes.

\begin{figure}[htbp]
\centering
\includegraphics[width=0.82\linewidth]{diagramme_mld}
\caption{Modele logique de donnees relationnel du systeme Z\"umm.}
\label{fig:mld}
\end{figure}

%------------------------------------------------------------------------------
\section{Maquettes de l'interface (wireframes)}
%------------------------------------------------------------------------------
Les maquettes basse-fidelite suivantes fixent l'ergonomie des ecrans cles avant
le developpement. Elles respectent la charte visuelle et la simplicite exigee.

\begin{figure}[htbp]
\centering
\includegraphics[width=0.9\linewidth]{wireframe_calendrier}
\caption{Maquette : calendrier / matrice agents x ruches avec pop-up de visite.}
\label{fig:wf-calendrier}
\end{figure}

\begin{figure}[htbp]
\centering
\includegraphics[width=0.9\linewidth]{wireframe_tdb_production}
\caption{Maquette : tableau de bord de production avec seuil parametrable.}
\label{fig:wf-production}
\end{figure}

\begin{figure}[htbp]
\centering
\includegraphics[width=0.95\linewidth]{wireframe_rapport_visite}
\caption{Maquette : formulaire de rapport de visite (avec mode hors-ligne et photos).}
\label{fig:wf-rapport}
\end{figure}

\begin{figure}[htbp]
\centering
\includegraphics[width=0.85\linewidth]{wireframe_ruche}
\caption{Maquette : fiche d'une ruche et composition (corps, hausses, cadres).}
\label{fig:wf-ruche}
\end{figure}

%------------------------------------------------------------------------------
\section{Matrice des droits d'acces (RBAC)}
%------------------------------------------------------------------------------
La matrice precise, pour chaque fonction, le droit accorde a chaque profil.
\textbf{Legende :} CRUD = creation / lecture / mise a jour / suppression ;
R = lecture ; X = execution ; A = approbation ; --- = aucun acces.

\begin{longtable}{>{\bfseries}p{3.6cm} C{1.3cm} C{1.3cm} C{1.3cm} C{1.3cm} C{1.3cm} C{1.3cm}}
\toprule
\textbf{Fonction} & \textbf{Prop.} & \textbf{Apic.} & \textbf{Sup.} & \textbf{Resp.} & \textbf{Admin} & \textbf{API} \\
\midrule
\endhead
Fermes et sites          & CRUD & R    & R    & R    & CRUD & --- \\
Ruches (corps/cadres)    & CRUD & R    & R    & R    & CRUD & --- \\
Agents                   & R    & ---  & R    & R    & CRUD & --- \\
Planifier une visite     & R    & X    & X    & X    & X    & --- \\
Approuver le planning    & ---  & ---  & A    & ---  & X    & --- \\
Realiser visite/rapport  & R    & X    & X    & ---  & R    & --- \\
Calendrier agents x ruches & R  & R    & R    & R    & R    & --- \\
TB production            & R    & ---  & ---  & R    & R    & --- \\
TB alertes / traiter     & R    & ---  & ---  & X    & R    & --- \\
Synthese et ROI          & R    & ---  & ---  & R    & R    & --- \\
Parametrer seuils/unites & ---  & ---  & ---  & R    & CRUD & --- \\
Internationalisation     & ---  & ---  & ---  & ---  & CRUD & --- \\
Export CSV / TXT         & R    & R    & R    & R    & R    & --- \\
Sauvegarde / restauration & --- & ---  & ---  & ---  & X    & --- \\
API quantite de miel     & ---  & ---  & ---  & ---  & ---  & R \\
\bottomrule
\end{longtable}

\begin{encadre}
Cette matrice doit etre appliquee cote serveur (verification systematique des
droits dans la couche de controle), et non uniquement par le masquage des
elements d'interface.
\end{encadre}

%------------------------------------------------------------------------------
\section{Registre des risques}
%------------------------------------------------------------------------------
\begin{longtable}{>{\bfseries}p{4.6cm} C{1.8cm} C{1.6cm} p{5.4cm}}
\toprule
\textbf{Risque} & \textbf{Probab.} & \textbf{Impact} & \textbf{Mesure de reduction} \\
\midrule
\endhead
Retard de developpement       & Moyenne & Fort      & Prioriser le noyau CRUD ; jalons hebdomadaires (S5 a S12). \\
Portee trop large             & Moyenne & Fort      & Se limiter aux fonctions de priorite haute ; le reste en option. \\
Complexite i18n arabe (RTL)   & Moyenne & Moyen     & Test precoce de bout en bout FR/EN/AR, gestion \texttt{dir=rtl}. \\
Perte de donnees terrain      & Moyenne & Fort      & File de commandes hors-ligne + synchronisation + sauvegarde. \\
Failles applicatives (injection, XSS) & Moyenne & Fort & \texttt{PreparedStatement}, echappement JSP, revue de securite. \\
Fausses alertes (capteurs)    & Moyenne & Moyen     & Seuils parametrables + validation humaine par le rapport. \\
Indisponibilite d'OAuth Google & Faible & Moyen     & Repli d'authentification locale pour la demonstration. \\
Volumetrie des mesures        & Faible  & Moyen     & Archivage et agregation par periode (volet automatise). \\
Usage d'un generateur de code & Faible  & Tres fort & Conception et code faits maison, tracabilite et revue par pairs. \\
\bottomrule
\end{longtable}

%------------------------------------------------------------------------------
\section{Protection des donnees personnelles}
%------------------------------------------------------------------------------
Le systeme manipule des donnees personnelles (noms et contacts des fermiers et
des agents) et des donnees potentiellement sensibles (\textbf{geolocalisation}
precise des sites). Les principes suivants sont retenus, dans l'esprit du RGPD :
\begin{itemize}
  \item \textbf{Minimisation et finalite} : ne collecter que les donnees utiles au suivi apicole.
  \item \textbf{Securite} : chiffrement des echanges (SSL / X.509) et acces authentifies, deja prevus.
  \item \textbf{Controle d'acces} : application stricte de la matrice RBAC ci-dessus.
  \item \textbf{Droits des personnes} : consultation, rectification et \textbf{effacement} des donnees, \textbf{portabilite} assuree par l'export CSV / TXT.
  \item \textbf{Conservation} : duree de retention definie pour les mesures et les visites, au-dela de laquelle les donnees sont archivees ou anonymisees.
  \item \textbf{Tracabilite} : journalisation des acces et des actions sensibles (approbation de planning, cloture de site).
\end{itemize}

%------------------------------------------------------------------------------
\section{Strategie de tests}
%------------------------------------------------------------------------------
La strategie complete le plan de tests fonctionnels (chapitre~10) par les
niveaux unitaire et d'integration. La conception faiblement couplee (interfaces
\texttt{Repository}, services segreges) rend les regles metier testables en
isolation, par injection de doubles (mocks).

\begin{longtable}{>{\bfseries}p{2.8cm} p{4.4cm} p{4.4cm} C{2.2cm}}
\toprule
\textbf{Niveau} & \textbf{Portee} & \textbf{Exemples} & \textbf{Outil} \\
\midrule
\endhead
Unitaire      & Regles de gestion isolees          & Contraintes cadres/hausses, calcul du miel (Composite), ROI, seuils & JUnit 5, Mockito \\
Integration   & Acces aux donnees (Repository/JDBC) & CRUD persiste, requetes de calcul                                   & JUnit + base H2 \\
Systeme       & Cas fonctionnels de bout en bout    & Cas TC01 a TC16 du chapitre~10                                      & JUnit, Selenium \\
Acceptation   & Validation par la demonstration     & Parcours agent, fermier, responsable                                & Soutenance \\
\bottomrule
\end{longtable}

\begin{itemize}
  \item \textbf{Couverture visee} : au moins 70\,\% de la couche metier (regles de gestion), mesuree par \texttt{JaCoCo}.
  \item \textbf{Tracabilite} : chaque exigence est reliee a au moins un cas de test (matrice du chapitre~10), completee par les tests unitaires correspondants.
  \item \textbf{Automatisation} : execution des tests a chaque build (Maven), pre-requis a la livraison.
\end{itemize}

\begin{encadre}
\textbf{Synergie avec la conception :} l'inversion de dependance (DIP) permet de
substituer un \texttt{RucheRepository} en memoire aux implementations JDBC, ce
qui rend les tests rapides, deterministes et independants de la base.
\end{encadre}
 dans cahier_des_charges.tex
%  (charte "Miel & noir" : miel, mielfonce, encadre, colonne C, longtable,
%   booktabs -- deja definis dans le preambule principal)
%==============================================================================

%==============================================================================
\chapter{Annexe - Complements : donnees, IHM, securite et tests}
\label{chap:complements}
%==============================================================================
Cette annexe regroupe des livrables complementaires renforcant le dossier :
le modele logique de donnees, les maquettes de l'interface, la matrice des
droits d'acces, le registre des risques, la conformite a la protection des
donnees et la strategie de tests.

%------------------------------------------------------------------------------
\section{Modele logique de donnees (MLD)}
%------------------------------------------------------------------------------
Le MLD traduit le dictionnaire de donnees (chapitre~5) en tables relationnelles
avec cles primaires, cles etrangeres et contraintes de gestion (\texttt{CHECK}
sur le nombre de hausses, la capacite en cadres et la productivite ; unicite du
couple compartiment/slot). Il distingue la relation d'\emph{emplacement}
(un site heberge des ruches) de la relation de \emph{propriete} (une ferme
possede des ruches), conformement a la regle autorisant un site a accueillir des
ruches de plusieurs fermes.

\begin{figure}[htbp]
\centering
\includegraphics[width=0.82\linewidth]{diagramme_mld}
\caption{Modele logique de donnees relationnel du systeme Z\"umm.}
\label{fig:mld}
\end{figure}

%------------------------------------------------------------------------------
\section{Maquettes de l'interface (wireframes)}
%------------------------------------------------------------------------------
Les maquettes basse-fidelite suivantes fixent l'ergonomie des ecrans cles avant
le developpement. Elles respectent la charte visuelle et la simplicite exigee.

\begin{figure}[htbp]
\centering
\includegraphics[width=0.9\linewidth]{wireframe_calendrier}
\caption{Maquette : calendrier / matrice agents x ruches avec pop-up de visite.}
\label{fig:wf-calendrier}
\end{figure}

\begin{figure}[htbp]
\centering
\includegraphics[width=0.9\linewidth]{wireframe_tdb_production}
\caption{Maquette : tableau de bord de production avec seuil parametrable.}
\label{fig:wf-production}
\end{figure}

\begin{figure}[htbp]
\centering
\includegraphics[width=0.95\linewidth]{wireframe_rapport_visite}
\caption{Maquette : formulaire de rapport de visite (avec mode hors-ligne et photos).}
\label{fig:wf-rapport}
\end{figure}

\begin{figure}[htbp]
\centering
\includegraphics[width=0.85\linewidth]{wireframe_ruche}
\caption{Maquette : fiche d'une ruche et composition (corps, hausses, cadres).}
\label{fig:wf-ruche}
\end{figure}

%------------------------------------------------------------------------------
\section{Matrice des droits d'acces (RBAC)}
%------------------------------------------------------------------------------
La matrice precise, pour chaque fonction, le droit accorde a chaque profil.
\textbf{Legende :} CRUD = creation / lecture / mise a jour / suppression ;
R = lecture ; X = execution ; A = approbation ; --- = aucun acces.

\begin{longtable}{>{\bfseries}p{3.6cm} C{1.3cm} C{1.3cm} C{1.3cm} C{1.3cm} C{1.3cm} C{1.3cm}}
\toprule
\textbf{Fonction} & \textbf{Prop.} & \textbf{Apic.} & \textbf{Sup.} & \textbf{Resp.} & \textbf{Admin} & \textbf{API} \\
\midrule
\endhead
Fermes et sites          & CRUD & R    & R    & R    & CRUD & --- \\
Ruches (corps/cadres)    & CRUD & R    & R    & R    & CRUD & --- \\
Agents                   & R    & ---  & R    & R    & CRUD & --- \\
Planifier une visite     & R    & X    & X    & X    & X    & --- \\
Approuver le planning    & ---  & ---  & A    & ---  & X    & --- \\
Realiser visite/rapport  & R    & X    & X    & ---  & R    & --- \\
Calendrier agents x ruches & R  & R    & R    & R    & R    & --- \\
TB production            & R    & ---  & ---  & R    & R    & --- \\
TB alertes / traiter     & R    & ---  & ---  & X    & R    & --- \\
Synthese et ROI          & R    & ---  & ---  & R    & R    & --- \\
Parametrer seuils/unites & ---  & ---  & ---  & R    & CRUD & --- \\
Internationalisation     & ---  & ---  & ---  & ---  & CRUD & --- \\
Export CSV / TXT         & R    & R    & R    & R    & R    & --- \\
Sauvegarde / restauration & --- & ---  & ---  & ---  & X    & --- \\
API quantite de miel     & ---  & ---  & ---  & ---  & ---  & R \\
\bottomrule
\end{longtable}

\begin{encadre}
Cette matrice doit etre appliquee cote serveur (verification systematique des
droits dans la couche de controle), et non uniquement par le masquage des
elements d'interface.
\end{encadre}

%------------------------------------------------------------------------------
\section{Registre des risques}
%------------------------------------------------------------------------------
\begin{longtable}{>{\bfseries}p{4.6cm} C{1.8cm} C{1.6cm} p{5.4cm}}
\toprule
\textbf{Risque} & \textbf{Probab.} & \textbf{Impact} & \textbf{Mesure de reduction} \\
\midrule
\endhead
Retard de developpement       & Moyenne & Fort      & Prioriser le noyau CRUD ; jalons hebdomadaires (S5 a S12). \\
Portee trop large             & Moyenne & Fort      & Se limiter aux fonctions de priorite haute ; le reste en option. \\
Complexite i18n arabe (RTL)   & Moyenne & Moyen     & Test precoce de bout en bout FR/EN/AR, gestion \texttt{dir=rtl}. \\
Perte de donnees terrain      & Moyenne & Fort      & File de commandes hors-ligne + synchronisation + sauvegarde. \\
Failles applicatives (injection, XSS) & Moyenne & Fort & \texttt{PreparedStatement}, echappement JSP, revue de securite. \\
Fausses alertes (capteurs)    & Moyenne & Moyen     & Seuils parametrables + validation humaine par le rapport. \\
Indisponibilite d'OAuth Google & Faible & Moyen     & Repli d'authentification locale pour la demonstration. \\
Volumetrie des mesures        & Faible  & Moyen     & Archivage et agregation par periode (volet automatise). \\
Usage d'un generateur de code & Faible  & Tres fort & Conception et code faits maison, tracabilite et revue par pairs. \\
\bottomrule
\end{longtable}

%------------------------------------------------------------------------------
\section{Protection des donnees personnelles}
%------------------------------------------------------------------------------
Le systeme manipule des donnees personnelles (noms et contacts des fermiers et
des agents) et des donnees potentiellement sensibles (\textbf{geolocalisation}
precise des sites). Les principes suivants sont retenus, dans l'esprit du RGPD :
\begin{itemize}
  \item \textbf{Minimisation et finalite} : ne collecter que les donnees utiles au suivi apicole.
  \item \textbf{Securite} : chiffrement des echanges (SSL / X.509) et acces authentifies, deja prevus.
  \item \textbf{Controle d'acces} : application stricte de la matrice RBAC ci-dessus.
  \item \textbf{Droits des personnes} : consultation, rectification et \textbf{effacement} des donnees, \textbf{portabilite} assuree par l'export CSV / TXT.
  \item \textbf{Conservation} : duree de retention definie pour les mesures et les visites, au-dela de laquelle les donnees sont archivees ou anonymisees.
  \item \textbf{Tracabilite} : journalisation des acces et des actions sensibles (approbation de planning, cloture de site).
\end{itemize}

%------------------------------------------------------------------------------
\section{Strategie de tests}
%------------------------------------------------------------------------------
La strategie complete le plan de tests fonctionnels (chapitre~10) par les
niveaux unitaire et d'integration. La conception faiblement couplee (interfaces
\texttt{Repository}, services segreges) rend les regles metier testables en
isolation, par injection de doubles (mocks).

\begin{longtable}{>{\bfseries}p{2.8cm} p{4.4cm} p{4.4cm} C{2.2cm}}
\toprule
\textbf{Niveau} & \textbf{Portee} & \textbf{Exemples} & \textbf{Outil} \\
\midrule
\endhead
Unitaire      & Regles de gestion isolees          & Contraintes cadres/hausses, calcul du miel (Composite), ROI, seuils & JUnit 5, Mockito \\
Integration   & Acces aux donnees (Repository/JDBC) & CRUD persiste, requetes de calcul                                   & JUnit + base H2 \\
Systeme       & Cas fonctionnels de bout en bout    & Cas TC01 a TC16 du chapitre~10                                      & JUnit, Selenium \\
Acceptation   & Validation par la demonstration     & Parcours agent, fermier, responsable                                & Soutenance \\
\bottomrule
\end{longtable}

\begin{itemize}
  \item \textbf{Couverture visee} : au moins 70\,\% de la couche metier (regles de gestion), mesuree par \texttt{JaCoCo}.
  \item \textbf{Tracabilite} : chaque exigence est reliee a au moins un cas de test (matrice du chapitre~10), completee par les tests unitaires correspondants.
  \item \textbf{Automatisation} : execution des tests a chaque build (Maven), pre-requis a la livraison.
\end{itemize}

\begin{encadre}
\textbf{Synergie avec la conception :} l'inversion de dependance (DIP) permet de
substituer un \texttt{RucheRepository} en memoire aux implementations JDBC, ce
qui rend les tests rapides, deterministes et independants de la base.
\end{encadre}
 dans cahier_des_charges.tex
%  (charte "Miel & noir" : miel, mielfonce, encadre, colonne C, longtable,
%   booktabs -- deja definis dans le preambule principal)
%==============================================================================

%==============================================================================
\chapter{Annexe - Complements : donnees, IHM, securite et tests}
\label{chap:complements}
%==============================================================================
Cette annexe regroupe des livrables complementaires renforcant le dossier :
le modele logique de donnees, les maquettes de l'interface, la matrice des
droits d'acces, le registre des risques, la conformite a la protection des
donnees et la strategie de tests.

%------------------------------------------------------------------------------
\section{Modele logique de donnees (MLD)}
%------------------------------------------------------------------------------
Le MLD traduit le dictionnaire de donnees (chapitre~5) en tables relationnelles
avec cles primaires, cles etrangeres et contraintes de gestion (\texttt{CHECK}
sur le nombre de hausses, la capacite en cadres et la productivite ; unicite du
couple compartiment/slot). Il distingue la relation d'\emph{emplacement}
(un site heberge des ruches) de la relation de \emph{propriete} (une ferme
possede des ruches), conformement a la regle autorisant un site a accueillir des
ruches de plusieurs fermes.

\begin{figure}[htbp]
\centering
\includegraphics[width=0.82\linewidth]{diagramme_mld}
\caption{Modele logique de donnees relationnel du systeme Z\"umm.}
\label{fig:mld}
\end{figure}

%------------------------------------------------------------------------------
\section{Maquettes de l'interface (wireframes)}
%------------------------------------------------------------------------------
Les maquettes basse-fidelite suivantes fixent l'ergonomie des ecrans cles avant
le developpement. Elles respectent la charte visuelle et la simplicite exigee.

\begin{figure}[htbp]
\centering
\includegraphics[width=0.9\linewidth]{wireframe_calendrier}
\caption{Maquette : calendrier / matrice agents x ruches avec pop-up de visite.}
\label{fig:wf-calendrier}
\end{figure}

\begin{figure}[htbp]
\centering
\includegraphics[width=0.9\linewidth]{wireframe_tdb_production}
\caption{Maquette : tableau de bord de production avec seuil parametrable.}
\label{fig:wf-production}
\end{figure}

\begin{figure}[htbp]
\centering
\includegraphics[width=0.95\linewidth]{wireframe_rapport_visite}
\caption{Maquette : formulaire de rapport de visite (avec mode hors-ligne et photos).}
\label{fig:wf-rapport}
\end{figure}

\begin{figure}[htbp]
\centering
\includegraphics[width=0.85\linewidth]{wireframe_ruche}
\caption{Maquette : fiche d'une ruche et composition (corps, hausses, cadres).}
\label{fig:wf-ruche}
\end{figure}

%------------------------------------------------------------------------------
\section{Matrice des droits d'acces (RBAC)}
%------------------------------------------------------------------------------
La matrice precise, pour chaque fonction, le droit accorde a chaque profil.
\textbf{Legende :} CRUD = creation / lecture / mise a jour / suppression ;
R = lecture ; X = execution ; A = approbation ; --- = aucun acces.

\begin{longtable}{>{\bfseries}p{3.6cm} C{1.3cm} C{1.3cm} C{1.3cm} C{1.3cm} C{1.3cm} C{1.3cm}}
\toprule
\textbf{Fonction} & \textbf{Prop.} & \textbf{Apic.} & \textbf{Sup.} & \textbf{Resp.} & \textbf{Admin} & \textbf{API} \\
\midrule
\endhead
Fermes et sites          & CRUD & R    & R    & R    & CRUD & --- \\
Ruches (corps/cadres)    & CRUD & R    & R    & R    & CRUD & --- \\
Agents                   & R    & ---  & R    & R    & CRUD & --- \\
Planifier une visite     & R    & X    & X    & X    & X    & --- \\
Approuver le planning    & ---  & ---  & A    & ---  & X    & --- \\
Realiser visite/rapport  & R    & X    & X    & ---  & R    & --- \\
Calendrier agents x ruches & R  & R    & R    & R    & R    & --- \\
TB production            & R    & ---  & ---  & R    & R    & --- \\
TB alertes / traiter     & R    & ---  & ---  & X    & R    & --- \\
Synthese et ROI          & R    & ---  & ---  & R    & R    & --- \\
Parametrer seuils/unites & ---  & ---  & ---  & R    & CRUD & --- \\
Internationalisation     & ---  & ---  & ---  & ---  & CRUD & --- \\
Export CSV / TXT         & R    & R    & R    & R    & R    & --- \\
Sauvegarde / restauration & --- & ---  & ---  & ---  & X    & --- \\
API quantite de miel     & ---  & ---  & ---  & ---  & ---  & R \\
\bottomrule
\end{longtable}

\begin{encadre}
Cette matrice doit etre appliquee cote serveur (verification systematique des
droits dans la couche de controle), et non uniquement par le masquage des
elements d'interface.
\end{encadre}

%------------------------------------------------------------------------------
\section{Registre des risques}
%------------------------------------------------------------------------------
\begin{longtable}{>{\bfseries}p{4.6cm} C{1.8cm} C{1.6cm} p{5.4cm}}
\toprule
\textbf{Risque} & \textbf{Probab.} & \textbf{Impact} & \textbf{Mesure de reduction} \\
\midrule
\endhead
Retard de developpement       & Moyenne & Fort      & Prioriser le noyau CRUD ; jalons hebdomadaires (S5 a S12). \\
Portee trop large             & Moyenne & Fort      & Se limiter aux fonctions de priorite haute ; le reste en option. \\
Complexite i18n arabe (RTL)   & Moyenne & Moyen     & Test precoce de bout en bout FR/EN/AR, gestion \texttt{dir=rtl}. \\
Perte de donnees terrain      & Moyenne & Fort      & File de commandes hors-ligne + synchronisation + sauvegarde. \\
Failles applicatives (injection, XSS) & Moyenne & Fort & \texttt{PreparedStatement}, echappement JSP, revue de securite. \\
Fausses alertes (capteurs)    & Moyenne & Moyen     & Seuils parametrables + validation humaine par le rapport. \\
Indisponibilite d'OAuth Google & Faible & Moyen     & Repli d'authentification locale pour la demonstration. \\
Volumetrie des mesures        & Faible  & Moyen     & Archivage et agregation par periode (volet automatise). \\
Usage d'un generateur de code & Faible  & Tres fort & Conception et code faits maison, tracabilite et revue par pairs. \\
\bottomrule
\end{longtable}

%------------------------------------------------------------------------------
\section{Protection des donnees personnelles}
%------------------------------------------------------------------------------
Le systeme manipule des donnees personnelles (noms et contacts des fermiers et
des agents) et des donnees potentiellement sensibles (\textbf{geolocalisation}
precise des sites). Les principes suivants sont retenus, dans l'esprit du RGPD :
\begin{itemize}
  \item \textbf{Minimisation et finalite} : ne collecter que les donnees utiles au suivi apicole.
  \item \textbf{Securite} : chiffrement des echanges (SSL / X.509) et acces authentifies, deja prevus.
  \item \textbf{Controle d'acces} : application stricte de la matrice RBAC ci-dessus.
  \item \textbf{Droits des personnes} : consultation, rectification et \textbf{effacement} des donnees, \textbf{portabilite} assuree par l'export CSV / TXT.
  \item \textbf{Conservation} : duree de retention definie pour les mesures et les visites, au-dela de laquelle les donnees sont archivees ou anonymisees.
  \item \textbf{Tracabilite} : journalisation des acces et des actions sensibles (approbation de planning, cloture de site).
\end{itemize}

%------------------------------------------------------------------------------
\section{Strategie de tests}
%------------------------------------------------------------------------------
La strategie complete le plan de tests fonctionnels (chapitre~10) par les
niveaux unitaire et d'integration. La conception faiblement couplee (interfaces
\texttt{Repository}, services segreges) rend les regles metier testables en
isolation, par injection de doubles (mocks).

\begin{longtable}{>{\bfseries}p{2.8cm} p{4.4cm} p{4.4cm} C{2.2cm}}
\toprule
\textbf{Niveau} & \textbf{Portee} & \textbf{Exemples} & \textbf{Outil} \\
\midrule
\endhead
Unitaire      & Regles de gestion isolees          & Contraintes cadres/hausses, calcul du miel (Composite), ROI, seuils & JUnit 5, Mockito \\
Integration   & Acces aux donnees (Repository/JDBC) & CRUD persiste, requetes de calcul                                   & JUnit + base H2 \\
Systeme       & Cas fonctionnels de bout en bout    & Cas TC01 a TC16 du chapitre~10                                      & JUnit, Selenium \\
Acceptation   & Validation par la demonstration     & Parcours agent, fermier, responsable                                & Soutenance \\
\bottomrule
\end{longtable}

\begin{itemize}
  \item \textbf{Couverture visee} : au moins 70\,\% de la couche metier (regles de gestion), mesuree par \texttt{JaCoCo}.
  \item \textbf{Tracabilite} : chaque exigence est reliee a au moins un cas de test (matrice du chapitre~10), completee par les tests unitaires correspondants.
  \item \textbf{Automatisation} : execution des tests a chaque build (Maven), pre-requis a la livraison.
\end{itemize}

\begin{encadre}
\textbf{Synergie avec la conception :} l'inversion de dependance (DIP) permet de
substituer un \texttt{RucheRepository} en memoire aux implementations JDBC, ce
qui rend les tests rapides, deterministes et independants de la base.
\end{encadre}
